The Schwartzberg score ($S = P / (2\sqrt{\pi A})$) is a perimeter-based
compactness metric on the scale $[1, \infty)$, where $S = 1$ represents a
perfect circle and larger values indicate greater non-compactness.
We prove that $S = 1/\sqrt{\text{PP}}$ exactly---meaning S and Polsby-Popper (PP)
carry identical geometric information with inverted scales and no information loss.
Despite this algebraic equivalence, S is preferred in legislative contexts
because its $[1, \infty)$ scale communicates compactness as ``distance from
circular'' more naturally than PP's inverted $[0,1]$ scale.
Colorado's redistricting constitution (Art.\ V, \S 44) mandates a
Schwartzberg-equivalent compactness criterion for its independent redistricting
commission, making S the legally operative metric for bisect results
submitted in Colorado proceedings.
We provide the full PP-to-S conversion table, a statutory survey of states
referencing S or equivalent perimeter-based criteria, and empirical S values
for NC ($k=14$), WI ($k=8$), TX ($k=38$), and CO ($k=8$) under all four
bisect structure algorithms (seed $= 42^\dagger$).
A district with PP $= 0.22$ has $S \approx 2.13$;
a district with PP $= 0.40$ has $S \approx 1.58$.
